\documentclass[10pt]{amsart}
\usepackage[a4paper,margin=23mm]{geometry}
\usepackage{amsmath,amssymb,amsthm,mathtools,hyperref,graphicx,xspace,letltxmacro,xpatch}
\usepackage[english]{babel}
\usepackage{xfrac}

\usepackage{mathrsfs}
\let\mathcal\mathscr

\providecommand{\eprintbio}[1]{\href{http://biorxiv.org/content/10.1101/#1}{\texttt{bioR{$\chi$}iv\string:\allowbreak#1}}}



\multlinegap0pt
\newenvironment{enumeratei}
{\bgroup\def\theenumi{\roman{enumi}}\def\theenumii{\arabic{enumii}}\begin{enumerate}}
{\end{enumerate}\egroup}
\def\mto{\mathchoice{\longmapsto}{\mapsto}{\mapsto}{\mapsto}}
\newcommand{\sbullet}{{\scriptscriptstyle\bullet}}
\newcommand{\csbullet}{{\raisebox{1pt}{$\sbullet$}}}
\newcommand{\Psfrac}[2]{(\sfrac{#1}{#2})}
\newcommand{\psfrac}[2]{\sfrac{(#1)}{#2}}
\newcommand{\spfrac}[2]{\sfrac{#1}{(#2)}}
\newcommand{\pspfrac}[2]{\sfrac{(#1)}{(#2)}}

\newcommand{\cB}{\mathcal{B}}
\newcommand{\cC}{\mathcal{C}}
\newcommand{\cI}{\mathcal{I}}
\newcommand{\cM}{\mathcal{M}}
\newcommand{\cN}{\mathcal{N}}
\newcommand{\cO}{\mathcal{O}}
\newcommand{\cS}{\mathcal{S}}
\newcommand{\cT}{\mathcal{T}}

\usepackage{tikz}
\usetikzlibrary{shapes}

\definecolor{green_mermoz}{rgb}{0,0.5,0}

\usepackage[noabbrev]{cleveref}
\creflabelformat{equation}{#2(#1)#3}
\crefname{thm}{Theorem}{Theorems}
\crefname{problem}{Problem}{Problems}
\crefname{prop}{Proposition}{Propositions}
\crefname{lem}{Lemma}{Lemmas}
\crefname{cor}{Corollary}{Corollaries}
\creflabelformat{problem}{#2(#1)#3}

\theoremstyle{plain}
\newtheorem{thm}{Theorem}[section]
\newtheorem{lem}[thm]{Lemma}
\newtheorem{prop}[thm]{Proposition}
\newtheorem{cor}[thm]{Corollary}
\newtheorem{defprop}[thm]{Definition-Proposition}

\theoremstyle{definition}
\newtheorem{definition}[thm]{Definition}
\newtheorem{remark}[thm]{Remark}

\newcommand{\loc}{\mathrm{loc}}
\newcommand{\comega}{\Omega^{\mathsf{c}}}

\renewcommand{\H}{\mathcal{H}_\eps}
\newcommand{\D}{\mathcal{D}_\eps}
\newcommand{\Je}{\mathcal{J}_\eps}

\newcommand{\N}{\mathbb{N}}
\newcommand{\RR}{\mathbb{R}}\let\R\RR
\let\eps\varepsilon
\let\ds\displaystyle

\newcommand{\abs}[1]{\left\lvert#1\right\rvert}
\newcommand{\Babs}[1]{\Bigl\lvert#1\Bigr\rvert}
\newcommand{\bbabs}[1]{\biggl\lvert#1\biggr\rvert}
\newcommand{\norm}[1]{\left\lVert#1\right\rVert}
\newcommand{\nalph}[1]{\left\lVert#1\right\rVert_\alpha}

\newcommand{\Ea}{\mathcal{E}^\alpha}
\newcommand{\Eaz}{\mathcal{E}^\alpha_0}
\newcommand{\QV}{G_\eps^V}
\newcommand{\dQVA}{dG_\eps^{V_1}}
\newcommand{\dQVB}{dG_\eps^{V_2}}
\newcommand{\dQVe}{dG_\eps^{V_\eps}}
\newcommand{\bz}{\overline{z}}

\newcommand{\Cgamma}{L_K(\gamma)}
\newcommand{\RK}{R_K}
\newcommand{\GK}{G_K}

\newcommand{\Gammam}{\Gamma^m_{\eps}}
\newcommand{\GammaI}{\Gamma^I_{\eps}}

\DeclareMathOperator{\Id}{Id}
\DeclareMathOperator{\lep}{\lambda_\varepsilon}
\DeclareMathOperator*{\argmin}{argmin}

\newcommand{\dint}{\displaystyle \int}
\newcommand{\vphia}{\varphi_\alpha}

\newcommand{\iep}[1]{I_\varepsilon (U_\varepsilon)(#1)}

\usepackage{xfrac,xcolor,graphicx}\graphicspath{{./figures/}}\providecommand{\skpt}{\par\smallskip\noindent}\providecommand{\pointir}{.\ }\providecommand{\eg}{e.g.}\providecommand{\ie}{i.e.}\providecommand{\resp}{resp.}\providecommand{\cf}{cf.}\providecommand{\etc}{etc.}\let\Subsection\subsection\let\Subsubsection\subsubsection
